\section{Reporting and ethics}
\label{sec:reporting-ethics}

The practical proposal is a reporting distinction. Participant studies should report recruitment, eligibility, consent, demographics where relevant, compensation, exclusion, and statistical treatment of responses. Expert-evaluation studies should report evaluator qualifications, evaluation task, independence, access to context, adjudication procedure, disagreement handling, and any relationship between evaluator and author.

The institutional claim needs care. Local research-ethics boards decide the scope of their own review, and the paper should not pretend to settle law or policy by definition. The argument is methodological and conceptual: ethics review should not be triggered by the word \enquote{judgment} alone. The relevant question is whether the project studies persons as persons, or uses professional evaluation to assess linguistic materials.

This distinction can improve ethics practice rather than weaken it. It routes real participant studies to participant protections and routes expert-evaluation studies to transparency, independence, conflict-management, and measurement-quality standards.
